\documentclass[11pt,a4paper]{article}
% Preâmbulo comum. Não duplicar pacotes nos módulos.
\usepackage[utf8]{inputenc}
\usepackage[T1]{fontenc}
\usepackage[brazil]{babel}
\usepackage{amsmath,amssymb,amsthm,mathtools}
\usepackage{geometry}
\geometry{a4paper,margin=2.3cm}
\usepackage{microtype}
\usepackage{enumitem}
\usepackage{graphicx}
\usepackage{xcolor}
\usepackage{booktabs}
\usepackage{hyperref}
\hypersetup{colorlinks=true,linkcolor=blue!40!black,urlcolor=blue!50!black}

\setcounter{tocdepth}{2}
\numberwithin{equation}{section}

\theoremstyle{plain}
\newtheorem{teorema}{Teorema}[section]
\newtheorem{proposicao}[teorema]{Proposição}
\newtheorem{lema}[teorema]{Lema}
\newtheorem{corolario}[teorema]{Corolário}

\theoremstyle{definition}
\newtheorem{definicao}[teorema]{Definição}
\newtheorem{exemplo}[teorema]{Exemplo}
\newtheorem{exercicio}{Exercício}[section]

\theoremstyle{remark}
\newtheorem{observacao}[teorema]{Observação}
\newtheorem{contraexemplo}[teorema]{Contraexemplo}
\newtheorem{intuicao}[teorema]{Intuição}

\newcommand{\R}{\mathbb{R}}
\newcommand{\N}{\mathbb{N}}
\newcommand{\Z}{\mathbb{Z}}
\newcommand{\Q}{\mathbb{Q}}
\DeclareRobustCommand{\eps}{\varepsilon}
\DeclareRobustCommand{\vet}[1]{\vec{#1}}
\DeclareRobustCommand{\dif}{\mathrm{d}}
\DeclareRobustCommand{\abs}[1]{\left\lvert #1 \right\rvert}
\DeclareRobustCommand{\norm}[1]{\left\lVert #1 \right\rVert}
\DeclareMathOperator{\grad}{grad}
\DeclareMathOperator{\Div}{div}
\DeclareMathOperator{\rot}{rot}

\newcommand{\base}[1]{\par\medskip\noindent\textbf{Base para relembrar.} #1\par\medskip}
\newcommand{\cuidado}[1]{\par\medskip\noindent\textbf{Quando a hipótese falha.} #1\par\medskip}

\newcommand{\figura}[2]{%
  \begin{figure}[htbp]
  \centering
  \IfFileExists{#1}{%
    \includegraphics[width=0.84\textwidth]{#1}%
  }{%
    \fbox{\parbox{0.8\textwidth}{\centering Figura ainda não gerada:\\ \texttt{#1}}}%
  }
  \caption{#2}
  \end{figure}%
}

\newenvironment{esqueleto}{%
  \par\medskip\noindent\textbf{Esqueleto da demonstração.}\par\nopagebreak
  \begin{enumerate}[leftmargin=1.4em]
}{\end{enumerate}\medskip}

\title{Caderno de prova: Green, Stokes e o campo conservativo}
\author{}
\date{}
\begin{document}
\maketitle
\tableofcontents

\noindent
Prove cada item desde as definições. Pode usar o teorema fundamental do cálculo, a regra da cadeia, o teorema de Clairaut, a troca da ordem de integração de uma função contínua num retângulo e a mudança polar no plano. Um item pode usar um item anterior. Não cite um teorema de cálculo vetorial que não esteja provado nesta folha.

\begin{exercicio}
Seja $U\subset\R^n$ aberto e seja $\phi\colon U\to\R$ de classe $C^1$. Seja $C$ uma curva de classe $C^1$ por pedaços contida em $U$, com extremidade inicial $A$ e extremidade final $B$. Prove que
\[
\int_C\nabla\phi\cdot\dif\vet{r}=\phi(B)-\phi(A).
\]
Em particular, a integral depende só dos extremos. Se $C$ é fechada, a integral é zero.
\end{exercicio}
\vspace{8cm}

\begin{exercicio}
Seja $\phi$ de classe $C^2$ num aberto de $\R^3$ e seja $\vet{F}=\nabla\phi$. Prove que $\nabla\times\vet{F}=\vet{0}$, escrevendo as três componentes. No plano, com $\vet{F}=(P,Q)=\nabla\phi$, prove que $\partial_x Q-\partial_y P=0$.
\end{exercicio}
\vspace{8cm}

\begin{exercicio}
Seja
\[
\vet{F}
=\Biggl(-\frac{y}{x^2+y^2},\;\frac{x}{x^2+y^2}\Biggr)
\]
definido em $\R^2\setminus\{\vet{0}\}$. Prove que $\partial_x Q-\partial_y P=0$ em todo ponto do domínio, calculando as parciais. Prove que a circulação de $\vet{F}$ na circunferência $x^2+y^2=1$, orientada no sentido anti-horário, vale $2\pi$. Conclua que $\vet{F}$ não é conservativo em $\R^2\setminus\{\vet{0}\}$.
\end{exercicio}
\vspace{8cm}

\begin{exercicio}
Seja $D=[a,b]\times[c,d]$, com $a<b$ e $c<d$, e sejam $P$ e $Q$ de classe $C^1$ num aberto que contém $D$. Oriente $\partial D$ no sentido anti-horário. Prove que
\[
\oint_{\partial D} P\,\dif x+Q\,\dif y
=\iint_D\bigl(\partial_x Q-\partial_y P\bigr)\,\dif A.
\]
\end{exercicio}
\vspace{8cm}

\begin{exercicio}
Sejam $a>0$ e $b>0$. Usando o teorema de Green, prove que a região $x^2/a^2+y^2/b^2\le 1$ tem área $\pi ab$. A fronteira pode ser parametrizada por $x=a\cos t$, $y=b\sin t$, $t\in[0,2\pi]$.
\end{exercicio}
\vspace{8cm}

\begin{exercicio}
Seja $R>0$ e seja $\vet{F}=(x,y,z)$. Prove, pelo teorema da divergência, que o fluxo de $\vet{F}$ para fora da esfera $x^2+y^2+z^2=R^2$ vale $4\pi R^3$. Prove o caso do teorema da divergência que usar. O volume da bola, se entrar na conta, deve ser calculado.
\end{exercicio}
\vspace{8cm}

\end{document}
